\documentclass{beamer}
\input{tbook-deck}
\title{From squares to matrix inequalities}
\author{Haobo Ma\inst{1} \and Wenlin Zhang\inst{1,2}}
\institute{\inst{1}Chrono AI, Singapore \and \inst{2}National University of Singapore, Singapore}
\date{2026-11-15}
\begin{document}
% Manuscript: Title; version 1, DOI 10.5281/zenodo.23920583
\begin{frame}
\titlepage
\note{Based on version 1 of the paper, preserving recorded authors and affiliations. Begin with size two, then present the rectangular theorem, stacking mechanism and precise evidence scope.}
\end{frame}
\section{Question}
% Manuscript: Section 1, posed inequality
\begin{frame}{How far does a scalar inequality go?}
\[
2xy\leq x^2+y^2,\qquad (x-y)^2\geq0.
\]
\begin{itemize}
\item Replace numbers by Hermitian matrices; use $XY+YX$.
\item Choose $k$ eigenvalues, multiply, then sum all choices: is there still a bound?
\item Yes, even for the absolute value of that sum.
\end{itemize}
\note{Start with the familiar square difference. Matrix multiplication need not commute, so use $XY+YX$, not $2XY$. There is no factor $1/2$. Credit Malkoun's first question in MO 512579, revision 1.}
\end{frame}
% Manuscript: Section 2, example after Corollary 4
\begin{frame}{Positive inputs, negative determinant}
\[
X=\begin{pmatrix}1&0\\0&0\end{pmatrix},\quad
Y=\begin{pmatrix}1&1\\1&1\end{pmatrix}.
\]
\begin{align*}
H=XY+YX&=\begin{pmatrix}2&1\\1&0\end{pmatrix},\\
R=X^2+Y^2&=\begin{pmatrix}3&2\\2&2\end{pmatrix}.
\end{align*}
\[
\det H=-1,\qquad \det R=2.
\]
\note{Explain positive semidefiniteness using $|v_1|^2$ and $|v_1+v_2|^2$. A star denotes conjugate transpose. Since $H$ is real symmetric, its negative determinant means opposite-sign eigenvalues. The modulus bound $1\leq2$ carries more information than $-1\leq2$.}
\end{frame}
% Manuscript: Section 1, notation
\begin{frame}{From trace to determinant}
$\lambda(M)$ lists all eigenvalues, with repetitions.
\[
e_k(t_1,\ldots,t_n)=\sum_{|S|=k}\prod_{i\in S}t_i.
\]
\begin{align*}
e_0&=1, &e_1(\lambda(M))&=\operatorname{tr}M,\\
e_2(a,b,c)&=ab+ac+bc, &e_n(\lambda(M))&=\det M.
\end{align*}
\begin{block}{Our size-two example}
$|e_1(\lambda(H))|=2\leq5$,\quad $|e_2(\lambda(H))|=1\leq2$.
\end{block}
\note{Use pairwise products of three numbers before naming elementary symmetric functions. The first is the sum of diagonal entries; the last is the determinant. The empty-product convention makes $e_0=1$.}
\end{frame}
\section{Result}
% Manuscript: Theorem 1
\begin{frame}{Written main result: rectangular inputs}
$A,B\in\mathbb C^{p\times n}$, $p,n\geq1$; define
\[
H=A^*B+B^*A,\qquad R=A^*A+B^*B.
\]
\begin{bookresult}{rectangular-bound}
For every $S\subseteq\{1,\ldots,n\}$,
\[
|\det H[S,S]|\leq\det R[S,S].
\]
For $0\leq k\leq n$,
\begin{align*}
|e_k(\lambda(H))|&\leq\sum_{|S|=k}|\det H[S,S]|\\
&\leq e_k(\lambda(R)).
\end{align*}
\end{bookresult}
\note{Explain that $M[S,S]$ selects the same row and column indices. This is source Theorem 1. The inputs need not be positive semidefinite or even square; the products are still square of size $n$. The empty set gives $1\leq1$.}
\end{frame}
\section{Mechanism}
% Manuscript: Section 2, coefficient identity
\begin{frame}{Replace spectra by small determinants}
Expand $I+tM$ by columns, or factor the same polynomial:
\begin{align*}
\det(I+tM)&=\sum_{k=0}^n t^k\sum_{|S|=k}\det M[S,S],\\
&=\prod_{i=1}^n(1+t\lambda_i(M)).
\end{align*}
Compare coefficients of $t^k$, $0\leq k\leq n$:
\[
e_k(\lambda(M))=\sum_{|S|=k}\det M[S,S].
\]
\alert{Bound each principal minor, then sum.}
\note{Choosing identity columns off $S$ leaves the corresponding principal minor with positive sign. The second expansion follows from the characteristic polynomial, so no invertibility or distinct eigenvalues are required.}
\end{frame}
% Manuscript: Theorem 1, proof
\begin{frame}{Swap the stacks, preserve the Gram matrix}
Fix $|S|=k$; let $A_S,B_S$ be the selected columns.
\[
U=\begin{pmatrix}A_S\\B_S\end{pmatrix},\quad
V=\begin{pmatrix}B_S\\A_S\end{pmatrix}
\quad (2p\times k).
\]
\begin{align*}
U^*U=V^*V&=R[S,S],\\
U^*V&=H[S,S].
\end{align*}
A Gram matrix records pairwise column inner products.
\note{First take real $a,b$ and explain equal-length vectors $u=(a,b)$ and $v=(b,a)$, then multiply the blocks. We never assume $A,B$ commute. Equal Gram matrices, rather than a common eigenbasis, drive this step.}
\end{frame}
% Manuscript: Lemma 2
\begin{frame}{Where does Cauchy–Schwarz act?}
\begin{bookresult}{gram-bound}
$U,V\in\mathbb C^{m\times k}$, $m,k\geq0$:
\[
|\det(U^*V)|^2\leq\det(U^*U)\det(V^*V).
\]
\end{bookresult}
Put $u_T=\det U[T,:]$, $v_T=\det V[T,:]$.
\[
\det(U^*V)=\sum_{|T|=k}\overline{u_T}v_T.
\]
Apply Cauchy–Schwarz to $(u_T)_T,(v_T)_T$.
\note{Cauchy–Binet identifies the two squared list norms with $\det(U^*U)$ and $\det(V^*V)$. Thus ordinary vector Cauchy–Schwarz acts on lists of minors. If $k>m$, rank is insufficient; if $k=0$, use the empty minor.}
\end{frame}
% Manuscript: Theorem 1, proof
\begin{frame}{From a squared bound to a spectral bound}
Equal Gram matrices give
\[
|\det H[S,S]|^2\leq(\det R[S,S])^2.
\]
$R[S,S]$ is positive semidefinite: take nonnegative roots.
\begin{align*}
|e_k(\lambda(H))|&\leq\sum_{|S|=k}|\det H[S,S]|\\
&\leq\sum_{|S|=k}\det R[S,S]\\
&=e_k(\lambda(R)),\quad 0\leq k\leq n.
\end{align*}
\note{Identify the triangle inequality, then the termwise minor bound. Positivity is needed for $R[S,S]$, not for $H$. Every step is finite algebra and includes singular matrices directly.}
\end{frame}
\section{Scope}
% Manuscript: Corollary 3; Section 3, source coverage
\begin{frame}{Return to the question; separate the scopes}
\begin{bookresult}{hermitian-bound}
Arbitrary Hermitian $X,Y\in\mathbb C^{n\times n}$; $n\geq1$, $0\leq k\leq n$:
\[
|e_k(\lambda(XY+YX))|\leq e_k(\lambda(X^2+Y^2)).
\]
\end{bookresult}
\begin{bookresult}{mo-positive}
Positive semidefinite Hermitian inputs, $1\leq k\leq n$:
\[
e_k(\lambda(XY+YX))\leq e_k(\lambda(X^2+Y^2)).
\]
\end{bookresult}
\note{Take $A=X,B=Y,p=n$ and use $X^*=X,Y^*=Y$. Only the lower statement is the exact scope of the public Lean result. The upper written corollary is stronger; keep their evidence separate.}
\end{frame}
% Manuscript: Corollary 4
\begin{frame}{Polynomial bound and simultaneous equality}
\begin{bookresult}{determinant-bound}
With $A,B,H,R$ as before, for every $z\in\mathbb C$:
\[
|\det(I+zH)|\leq\det(I+|z|R).
\]
\[
|\det H|\leq\det R,\quad
\operatorname{tr}R-\operatorname{tr}H=\|A-B\|_F^2.
\]
All one-sided $e_k$ equalities, $1\leq k\leq n$, hold together
$\Longleftrightarrow A=B$.
\end{bookresult}
Necessity uses only $k=1$; absolute equalities differ.
\note{Bound the polynomial coefficients termwise and use the triangle inequality. The squared Frobenius norm sums squared entry moduli. With $A=-B\ne0$, every absolute-value equality holds even though $A\ne B$.}
\end{frame}
% Manuscript: Sections 1 and 3; bibliography
\begin{frame}{Meaning, sources and next steps}
\begin{itemize}
\item One mechanism controls trace, determinant and intermediate sums.
\item No commutation; no approximation for singular inputs.
\item The three-matrix Schur result requires pairwise commutation.
\end{itemize}
\small
Malkoun, MO 512579 revision 1 (2026).\\
J. K. Merikoski, A. Virtanen (1992),
\href{https://doi.org/10.21136/cmj.1992.128332}{compound matrices and symmetric spectral functions}.\\
Haobo Ma, Wenlin Zhang,
\href{https://doi.org/10.5281/zenodo.23920583}{version 1 paper}.\\
Shared requests: Gram bound; rectangular extension; polynomial and equality; commuting Schur.
\note{End the main story here. Contribution links are in the reading chapter. Credit cubesteak's posted proof of the different normalized three-matrix determinant problem, MO answer 514685, 26 August 2026. This talk does not settle all three-matrix questions.}
\end{frame}
\appendix
\section{Boundaries}
% Manuscript: Lemma 2 and proofs of Theorem 1, Corollary 4
\begin{frame}{Empty minors, rank and zero parameters}
\begin{itemize}
\item $k=0$: empty minors and products are $1$.
\item Gram lemma, $k>m$: all three determinants vanish.
\item Stacks, $k>2p$: rank makes the relevant minors zero.
\item $z=0$: the polynomial bound is $1=1$.
\item Repeated eigenvalues count with multiplicity; no limits.
\end{itemize}
\note{These boundaries follow directly from finite determinant identities, rather than an appeal to continuity. Finite examples do not replace this all-parameter argument.}
\end{frame}
\section{Three matrices}
% Manuscript: Corollary 5
\begin{frame}{Commuting Schur siblings}
Pairwise commuting PSD Hermitian $A,B,C$, $n\geq1$:
\[
S=A+B+C,\quad Q=A^2+B^2+C^2,
\]
\[
F=6(A^3+B^3+C^3)+S^3.
\]
\begin{bookresult}{commuting-schur}
\[
\det F\geq5^n\det S\det Q.
\]
For $0\leq k\leq n$,
\begin{align*}
e_k(\sigma(F))&\geq e_k(\sigma(5SQ)),\\
e_k(\lambda(2F))&\geq e_k(\lambda(5SQ+5QS)).
\end{align*}
\end{bookresult}
\note{$\sigma$ lists singular values with multiplicity. This is source Corollary 5. All three clauses require commutation. It is independent of the rectangular theorem and does not establish the noncommuting versions.}
\end{frame}
% Manuscript: Corollary 5, proof
\begin{frame}{The scalar calculation after diagonalization}
In a common eigenbasis, sort $a\geq b\geq c\geq0$ at each position.
\[
s=a+b+c,\quad q=a^2+b^2+c^2,
\]
\[
f=6(a^3+b^3+c^3)+s^3.
\]
\begin{align*}
f-5sq&=2\bigl((a-b)^2(a+b-c)\\
&\hspace{2em}+c(a-c)(b-c)\bigr)\geq0.
\end{align*}
Multiply entries, then sum; use $SQ=QS$ for the last clause.

\note{Commutation preserves eigenspaces and supplies a common unitary eigenbasis. At each position we can sort the three scalars independently because the expression is symmetric. All factors are nonnegative; no division by a zero entry is needed.}
\end{frame}
% Manuscript: Source coverage and verification
\FormalVerification
\end{document}
